\section{Materials and methods}

\subsection{Fast reading and reduction to \SI{1}{\kilo\hertz}}
The file is split into 1024 blocks whose lines are counted in parallel to build a row index,
so any time window can be read directly. The blocks are then reduced to the mean, minimum and
maximum of every 50 samples (\SI{50}{\kilo\hertz} $\rightarrow$ \SI{1}{\kilo\hertz}).
The signal has no content above about \SI{500}{\hertz}, so the \SI{1}{\milli\second} mean keeps
99.3--99.9\,\% of the peak height. The index and summary are cached for later opens.

\subsection{Preprocessing (Part A)}
Part A works on the \SI{1}{\kilo\hertz} means of each channel and has three steps:
\begin{enumerate}
  \item a median filter (3 samples), which removes single-sample spikes;
  \item a Savitzky--Golay filter (7 samples, order 3)~\cite{savitzky1964}, which smooths small
        fluctuations;
  \item subtraction of a slowly varying baseline. The baseline is fitted by asymmetric least
        squares (ALS)~\cite{eilers2005} on \SI{10}{\hertz} block means (smoothness
        $\lambda = 10^{5}$, asymmetry $p = 0.01$, 10 iterations) and interpolated back to
        \SI{1}{\kilo\hertz}. Samples above the fitted curve get a small weight, so the baseline
        follows the signal from below and is not pulled up by the peaks.
\end{enumerate}
The ALS equations form a symmetric pentadiagonal system, which is solved with a banded
Cholesky solver, and the two channels are processed in parallel threads. On all recordings
compared, this gives the same pieces and labels as the earlier sparse-matrix implementation.
If the median ratio \texttt{adc1}/\texttt{adc2} at the 50 largest \texttt{adc2} peaks is below
0.1, \texttt{adc1} is not used for the two-channel check.

\subsection{Noise level and threshold}
\label{sec:threshold}
The noise level $\sigma$ of each channel is estimated from the cleaned signal itself. The signal
is divided into \SI{200}{\milli\second} blocks, blocks whose standard deviation is below half the
median block standard deviation are dropped as idle (no flow), and $\sigma$ is the root mean
square of the standard deviations of the quietest 10\,\% of the remaining blocks.
Part B detects peaks whose prominence exceeds $k\sigma$ with $k = 10$.
All settings are the same for every recording: $k$, the \SI{1}{\second} grouping gap and the
no-dye level of 68.9 counts. Only $\sigma$ differs between recordings, and it is measured from
each signal and not chosen by concentration.

\subsection{Detection and classification (Part B)}
Peaks with prominence above $10\sigma$ are detected on both channels. Peaks less than
\SI{1}{\second} apart are grouped into one piece. Each piece is labelled from its largest peak
(Table~\ref{tab:labels}).

\begin{table}[h]
  \centering
  \caption{Labelling rules.}
  \label{tab:labels}
  \begin{tabular}{ll}
    \toprule
    Condition & Label \\
    \midrule
    above the no-dye level (68.9 counts) & \texttt{dyed} \\
    below, confirmed by \texttt{adc1} & \texttt{solid\_no\_dye} \\
    below, not confirmed & \texttt{water} \\
    below, two-channel check not usable & \texttt{undecided} \\
    \bottomrule
  \end{tabular}
\end{table}

The two-channel check requires an \texttt{adc1} peak within \SI{\pm 20}{\milli\second} of the
top \texttt{adc2} peak (shifted by the measured channel lag) and at least 5\,\% as high.
The chance rate is estimated by shifting \texttt{adc1} by 2--\SI{20}{\second}; above 20\,\%
the check is disabled.

\subsection{Amount of tissue}
The amount is the total area (counts\,$\times$\,ms) of the \texttt{dyed} and
\texttt{solid\_no\_dye} pieces. No conversion to millilitres is made yet, since the area per
millilitre depends on the dye concentration and on re-suction.

\subsection{Visualizer}
% TODO: describe the Streamlit app (overview with draggable window, detail view, heatmap, events table, LTTB).
